\section{Conditional regret and the first crossing}\label{sec:regret}

Fix a seed table satisfying the hypothesis of Theorem~\ref{thm:refinement},
and fix the comparison tree from Section~\ref{sec:comparison}.
We will bound the error accumulated at its nodes over the refinement indices.
The estimates may lie below true cost when $a>1$. To keep such
underestimation from cancelling an error elsewhere, define, for the analysis
only,
\[
 \overline d_t(I,q)=\max\{c(I,q),d_t(I,q)\}.
\]
At a nonterminal comparison node use the notation
$c,b,r,e,q_i,m_i,W,G,K$ of Section~\ref{sec:comparison}, and set
\begin{equation}\label{eq:node-regret-definition}
 \mathcal R_t(I,q)
 =\overline d_t(I,q)-Le-\sum_i\overline d_t(I_i,q_i).
\end{equation}
These quantities telescope over the fixed comparison tree. We first bound
their sum in time at an arbitrary node, then improve that bound on a node
whose comparison action belongs to a computed band.

\subsection{A telescoping bound at every comparison node}

\begin{lemma}\label{lem:generic-regret}
At every nonterminal comparison node,
\begin{equation}\label{eq:generic-regret}
 \sum_{t=1}^S\mathcal R_t
 \le \sum_i\bigl(\overline d_0-\overline d_S\bigr)(I_i,q_i)
 \le \sum_i(m_i-|z_i|)
 \le W-|z(I,r)|,
\end{equation}
where $z_i=z(I_i,q_i)$. In particular, the bound is at most $4Fc$.
If the node is not bad and $|\widehat v|>1-\tau$, it is at most
$(H^{-1/4}+4\tau F)c$.
\end{lemma}

\begin{proof}
The Bellman option and the comparison action give
\[
 d_t(I,q)\le Le+\sum_i d_{t-1}(I_i,q_i).
\]
The right side with all child values replaced by their clipped values is
at least $c$, by Equation~\ref{eq:comparison-split}. It therefore also bounds
$\overline d_t(I,q)$. Subtracting the same expression at index $t$ and
summing proves the first inequality in Equation~\ref{eq:generic-regret}.

For the second inequality, consider a child separately. If $d_0\le c_i$,
monotonicity makes both clipped endpoint values equal to $c_i$, so their
difference is zero, while $m_i\ge |z_i|$. If $d_0>c_i$, then
$\overline d_0=d_0\le m_i$ and $\overline d_S\ge |z_i|$. This proves the
required difference bound in both cases; it does not require
$\overline d_0\le m_i$ in the first case. Finally,
$\sum_i|z_i|\ge|z(I,r)|$ and $W\le4Fc$.

For the last assertion, $|\widehat v|=|z(I,r)|/Z>1-\tau$ implies
\[
 W-|z(I,r)|\le G+\tau W
 \le (H^{-1/4}+4\tau F)c,
\]
using $G\le bH^{-1/4}$ and $b\le c$ at a node that is not bad.
\end{proof}

\subsection{An entropy inequality with fixed block masses}

Exponential weighting and its potential analysis have a classical role in
online learning; see Freund and Schapire~\cite{FreundSchapire1997}.
Our predicted objective has the form of optimistic regularized leader,
as studied by Rakhlin and Sridharan~\cite[Section~2]{RakhlinSridharan2013}.
For our action vectors, each child receives a fixed total probability mass.
Cost decrements may differ greatly between children, but their variation
within a child is small.  The following deterministic lemma isolates the
reason that this weaker variation bound suffices.

\begin{lemma}[Monotone costs with fixed block masses]\label{lem:entropy}
Let $\mathcal Q$ be a nonempty compact convex set of probability vectors
in $\mathbb R^d$, containing a vector with every coordinate positive.
Partition the coordinates into blocks, and suppose that the total mass
in each block is the same for every vector in $\mathcal Q$.
For an integer $S\ge1$, let $f_0,\ldots,f_S\in[0,1]^d$ be
coordinatewise nonincreasing.
For $\eta>0$ and $0<\kappa<1/2$, set
\[
 \widehat g_t=-f_{t-1},\qquad
 \Delta_t=f_{t-1}-f_t,\qquad
 u_t=\widehat g_t+(1-2\kappa)\Delta_t.
\]
Suppose that the range of $\Delta_t$ within each block is at most
$\kappa/\eta$ for every $1\le t\le S$.  Let $\rho_t$ maximize
\[
 \mathcal H(\rho)+
 \eta\left\langle\sum_{j<t}u_j+\widehat g_t,\rho\right\rangle
 \quad\text{over }\rho\in\mathcal Q,
 \qquad \mathcal H(\rho)=-\sum_k\rho_k\ln\rho_k,
\]
with $0\ln0=0$.  For every $\rho^*\in\mathcal Q$ and every integer
$0\le s\le S$,
\begin{equation}\label{eq:entropy-regret}
\begin{split}
 \sum_{t=1}^s
 \left\langle f_t+\kappa\Delta_t,\rho_t\right\rangle
 \le {}&
 \sum_{t=1}^s
 \left\langle f_t,\rho^*\right\rangle
 +\frac{\ln d}{\eta}+2\kappa.
\end{split}
\end{equation}
For any $\zeta\ge0$, replacing each $\rho_t$ on the left by a vector
within $\zeta$ in $\ell_1$ distance adds at most $s\zeta$ to the bound.
\end{lemma}

\begin{proof}
The monotone cost vectors satisfy, coordinatewise,
\begin{equation}\label{eq:increment-properties}
 \Delta_t\ge0,
 \qquad
 \sum_{t=1}^S\Delta_t
 =f_0-f_S\le1.
\end{equation}
Define the potential
\[
 \Psi(v)=\max_{\rho\in\mathcal Q}
 \{\mathcal H(\rho)+\eta\langle v,\rho\rangle\}.
\]
Compactness gives a maximizer, and strict concavity of entropy makes it
unique.  It is positive: moving a small
distance $u$ from a vector with a zero coordinate toward the positive feasible
vector increases entropy by a positive multiple of $u\ln(1/u)+O(u)$,
whereas the linear term changes by only $O(u)$. At a positive maximizer
$p$, first-order optimality along feasible segments gives
\[
 \langle\nabla\mathcal H(p)+\eta v,\rho-p\rangle\le0
 \qquad(\rho\in\mathcal Q).
\]
Since
$\mathcal H(\rho)-\mathcal H(p)
 =\langle\nabla\mathcal H(p),\rho-p\rangle-\KL(\rho\Vert p)$,
we obtain the precise penalty
\begin{equation}\label{eq:entropy-penalty}
 \mathcal H(\rho)+\eta\langle v,\rho\rangle
 \le \Psi(v)-\KL(\rho\Vert p).
\end{equation}
Here $\KL(\rho\Vert p)=\sum_k\rho_k\ln(\rho_k/p_k)$, with zero terms
interpreted as zero. The penalty therefore holds on $\mathcal Q$ itself,
even when its feasible directions are restricted.

Put $v_{t-1}=\sum_{j<t}u_j$.
Evaluating the old objective at the maximizer $\rho_t$ of the new one
gives
\begin{equation}\label{eq:predicted-potential-increment}
 \Psi(v_{t-1}+\widehat g_t)-\Psi(v_{t-1})
 \le\eta\langle\widehat g_t,\rho_t\rangle.
\end{equation}
This inequality remains valid when the increment is negative.
To add the remaining vector $(1-2\kappa)\Delta_t$, let $\mu$ be the
vector whose coordinates are the minimum of $\Delta_t$ in their
block, and write $v'=\Delta_t-\mu$. Then
$0\le\eta v'_k\le\kappa$ and $\mu\ge0$.
Moreover $\langle\mu,\rho\rangle$ is the same for every
$\rho\in\mathcal Q$, because the block masses are fixed.
Thus the possibly large block minima contribute exactly the same linear
increment at every feasible vector; only the small within-block remainder
needs an exponential estimate.
Equation~\ref{eq:entropy-penalty} thus bounds the remaining potential
increment by
\begin{equation}\label{eq:innovation-potential-increment}
 \eta(1-2\kappa)\langle\mu,\rho_t\rangle
 +\ln\sum_k\rho_{t,k}
       \exp\bigl(\eta(1-2\kappa)v'_k\bigr).
\end{equation}
To justify the logarithmic term, extend the supremum with its relative
entropy penalty to the full probability simplex and normalize the weights
$\rho_{t,k}\exp(\eta(1-2\kappa)v'_k)$.
The resulting objective equals the displayed logarithm minus relative
entropy to these normalized weights, whose nonnegativity proves the bound.

For $0\le x\le\kappa\le1/2$,
\begin{equation}\label{eq:entropy-exponential}
 e^{(1-2\kappa)x}\le1+(1-\kappa)x.
\end{equation}
Indeed the derivative of
$\ln(1+(1-\kappa)x)-(1-2\kappa)x$ is
\[
 \frac{\kappa-(1-2\kappa)(1-\kappa)x}{1+(1-\kappa)x}\ge0,
\]
and the expression vanishes at zero.
Apply Equation~\ref{eq:entropy-exponential} with
$x=\eta v'_k\le\kappa$, and then use
$\ln(1+y)\le y$. Because $\mu\ge0$,
Equation~\ref{eq:innovation-potential-increment} is at most
$\eta(1-\kappa)\langle\Delta_t,\rho_t\rangle$.
Combining the two increments yields
\[
 \Psi(v_t)-\Psi(v_{t-1})
 \le\eta\left\langle
      \widehat g_t+(1-\kappa)\Delta_t,\rho_t\right\rangle.
\]

Sum through $s$, use $\Psi(0)\le\ln d$, and evaluate the
last potential at $\rho^*$. Its entropy is nonnegative. Replacing the
coefficient $1-2\kappa$ in the accumulated comparator objective by $1$
loses at most
\[
 2\kappa\left\langle\sum_{t=1}^s\Delta_t,\rho^*\right\rangle
 \le2\kappa
\]
by Equation~\ref{eq:increment-properties}.  Finally,
\[
 \widehat g_t+\Delta_t=-f_t,\qquad
 \widehat g_t+(1-\kappa)\Delta_t=-(f_t+\kappa\Delta_t).
\]
Reversing the signs gives~\eqref{eq:entropy-regret}.
The cost vector $f_t+\kappa\Delta_t=\kappa f_{t-1}+(1-\kappa)f_t$
lies in $[0,1]^d$, so an $\ell_1$ error of at most $\zeta$ costs at most
$\zeta$ per term.
\end{proof}

Apply the lemma to an unskipped band $\mathcal C=\mathcal C_{b,r}$,
conditioning on the complete child tables, including the warmup and all
indices through $S$.  The band, the scales $h_i$, and its action hull
$\mathcal Q$ are then fixed.  Take the child coordinates as blocks and set
\[
 f_t(i,s)=d_t(I_i,s)/h_i.
\]
The action hull is compact and convex, and averaging its action vectors
gives a positive feasible vector.  Every vector has mass $h_i/h_\Sigma$
in child block $i$.  The refinement invariants give
$0\le d_S\le\cdots\le d_0\le w_0\le h_i$ on band coordinates, so the
normalized costs meet the monotonicity and unit bounds of the lemma.

Two states in a child block have endpoint-pair distance at most
$4\theta b/M$.  The spatial bound~\eqref{eq:refinement-lipschitz}, used
in both directions and with $d_t\le h_i$, bounds their cost difference
at one index by $\delta h_i+4L\theta b/M$; the stronger bound at index
zero also suffices.  Since $h_i\ge b/M$, the within-block range of
$\Delta_t$ is at most
\begin{equation}\label{eq:increment-range}
 2\delta+8L\theta\le3\delta\le\kappa/\eta.
\end{equation}
Here $8L\theta=1024\ell^{-39}\le\ell^{-10}$ and
$3\eta\delta=3\ell^{-5}\le\kappa$.
The lemma therefore applies to exactly the objective
\eqref{eq:online-objective}, with $d=d_{\mathcal C}$.
Conditioning on child tables leaves the parent group samples and online
draws with their prescribed laws, since their streams are separate from
descendant streams.  No independence among child tables is needed.

To express the conclusion in action costs, put
\[
 x_t(\alpha)=\sum_i d_t(I_i,s_i(\alpha)).
\]
The fixed coordinate weights give
\[
 \begin{aligned}
 h_\Sigma\langle f_t,\rho(\alpha)\rangle&=x_t(\alpha),\\
 h_\Sigma\langle f_t+\kappa\Delta_t,\rho(\alpha)\rangle
 &=x_t(\alpha)+\kappa\bigl(x_{t-1}(\alpha)-x_t(\alpha)\bigr).
 \end{aligned}
\]
If $\alpha_t$ is drawn from the computed mixture, its expected action
vector is that mixture's mean.  Lemma~\ref{lem:entropy}, with
$\zeta=H^{-3}$, thus gives every fixed band action $\alpha^*$ and prefix
$0\le s\le S$ the conditional bound
\begin{equation}\label{eq:band-action-regret}
\begin{split}
 \E\left[\sum_{t=1}^s
   \left\{x_t(\alpha_t)
       +\kappa\bigl(x_{t-1}(\alpha_t)-x_t(\alpha_t)\bigr)\right\}
       \,\middle|\,\text{child tables}\right]
 \le{}&\sum_{t=1}^s x_t(\alpha^*)\\
 &+h_\Sigma\left(\frac{\ln d_{\mathcal C}}{\eta}
                         +2\kappa+sH^{-3}\right).
\end{split}
\end{equation}
The next argument shows why the decrement term on the left is enough
to cover the clipping in the parent update.

\subsection{Clipping before and after the first crossing}

We now apply the entropy inequality at a comparison node that is not bad
and satisfies $|\widehat v|\le1-\tau$.
Lemma~\ref{lem:comparison-band} ensures that its comparison action
$\alpha^*$ belongs to the unskipped band $\mathcal C_{b,r}$ and that
all its guards pass. Write
\[
 A_t(\alpha)=Le+x_t(\alpha).
\]
The entropy bound controls these action costs. The following proposition
transfers that control to the clipped parent regret.

\begin{proposition}\label{prop:node-regret}
Condition on the complete child tables at such a comparison node.
Then
\begin{equation}\label{eq:node-regret}
\begin{split}
 \E\left[\sum_{t=1}^S\mathcal R_t
       \,\middle|\,\text{child tables}\right]
 \le {}&h_\Sigma\left(
       \frac{\ln d_{\mathcal C}}{\eta}
       +2\kappa+SH^{-3}\right)\\
       &+2\delta c+2SH^{-4}Fc.
\end{split}
\end{equation}
Here $h_\Sigma\le64Fb$ and
$\ln d_{\mathcal C}\le4(p_*+21)\ell$.
\end{proposition}

\begin{proof}
Let $t_\times$ be the first $t\in\{1,\ldots,S\}$ for which
$\min_{\alpha\in\mathcal C}A_t(\alpha)<c$, or put $t_\times=S+1$
if there is none. This index depends only on the conditioned child tables.
It is an analysis device, not a test performed by the algorithm.
We will use~\eqref{eq:band-action-regret} on the fixed prefix
$1\le t<t_\times$.

Let $E_t$ be the failure event for the group approximation associated
with this representative. By Lemma~\ref{lem:group-accuracy},
$\Pr(E_t\mid\text{child tables})\le H^{-4}$, and on its complement
Equation~\ref{eq:group-accuracy} holds for every action in the band.
For the online action $\alpha_t$ sampled for this band, we claim that
on the prefix
\begin{equation}\label{eq:prefix-clipping}
\begin{split}
 \overline d_t(I,q)
 \le {}&Le+x_t(\alpha_t)
       +\kappa\bigl(x_{t-1}(\alpha_t)-x_t(\alpha_t)\bigr)\\
       &+2Fc\,\ind_{E_t}.
\end{split}
\end{equation}
Every action cost $A_t(\alpha)$ is at least $c$ on this prefix, and the
online candidate satisfies $Y_t(q)\le A_t(\alpha_t)$.
If $\overline d_t(I,q)\le A_t(\alpha_t)$, the claim follows at once
from child monotonicity. Otherwise, the update
Equation~\ref{eq:refined-update} forces
\[
 \Pi_t(q)>(1+\delta)A_t(\alpha_t),
 \qquad d_t(I,q)=\frac{\Pi_t(q)}{1+\delta}.
\]
The earlier-index Bellman option also gives
$\Pi_t(q)\le Le+x_{t-1}(\alpha_t)$, so, writing
$\Delta x=x_{t-1}(\alpha_t)-x_t(\alpha_t)$,
\[
 \Delta x>\delta A_t(\alpha_t)\ge\delta c\ge\delta b.
\]
On $E_t^c$, Equation~\ref{eq:group-accuracy} and $e_*\le\kappa\delta$
now give the full excess calculation
\[
 \begin{split}
 d_t(I,q)-A_t(\alpha_t)
 &\le\frac{e_*b-\delta Le}{1+\delta}
 \le\frac{e_*b}{1+\delta}\\
 &\le\frac{\kappa\delta b}{1+\delta}
 \le\kappa\Delta x.
 \end{split}
\]
On $E_t$, the added term $2Fc$ suffices because
$\overline d_t(I,q)\le2Fc$. This proves
Equation~\ref{eq:prefix-clipping} without requiring independence
between the sampled action and the group event.

If $t_\times\le S$, choose an action attaining the crossing.
On $E_{t_\times}^c$, its group bound gives
\[
 \Pi_{t_\times}(q)
 \le Le+(1+\delta)x_{t_\times}(\alpha)+e_*b
 <c+\delta c+e_*b\le c+2\delta c.
\]
Thus the clipped regret at that index is at most $2\delta c$,
with an additional $2Fc$ allowed on failure. At every later index,
the same crossing action and child monotonicity imply
$\mathcal B_I(d_{t-1})(q)<c$. Hence $d_t(I,q)<c$,
$\overline d_t(I,q)=c$, and the regret is nonpositive by
Equation~\ref{eq:comparison-split}.

It remains to sum~\eqref{eq:prefix-clipping} on this prefix.
Use~\eqref{eq:band-action-regret} with $s=t_\times-1$ and the
comparison action $\alpha^*$.  Since
$x_t(\alpha^*)\le\sum_i\overline d_t(I_i,q_i)$, its comparator costs
are bounded by the child terms subtracted in each regret.  The expected
prefix regret is therefore at most
\[
 h_\Sigma\left(\frac{\ln d_{\mathcal C}}{\eta}
                  +2\kappa+SH^{-3}\right),
\]
apart from the explicit failure charges in~\eqref{eq:prefix-clipping}.
The prefix and possible crossing contain at most $S$ failure events,
each of conditional probability at most $H^{-4}$. Adding their expected
costs and the crossing charge proves Equation~\ref{eq:node-regret}.

Finally each child has at most $(21P)^2$ possible states within the wide
shift radius. Thus $d_{\mathcal C}\le M(21P)^2$, from which
$\ln d_{\mathcal C}\le4(p_*+21)\ell$ follows using
$\log_2 H\le\ell$. The bound on $h_\Sigma$ is one of the band guards.
\end{proof}

The expectation in Proposition~\ref{prop:node-regret} concerns one node
after fixing its child tables. Integrating it and summing over the fixed
comparison tree uses only linearity of expectation. No event asserting
that every group in the recursive computation succeeds is required.
